% conclusions.tex (국문 v13)
\section{결론}\label{sec:conclusions}
이 연구는 라벨, 기후 입력, 채점이 같은 관측 지점과 같은 해를 가리키는 지점-연도 설계로 활동층 두께 모델을 평가하였다. 기후 입력을 관측 연도와 맞추면 관측이 여러 해에 분포한 지역에서 공간 블록 오차가 감소하고, 모든 예측이 명시된 해에 대응한다. 지역 라벨로 재보정한 Stefan 모델에 gradient boosting 잔차 모델을 더한 물리 잔차 모델은 Stefan 모델의 오차를 레나 델타에서 \ndrefLeblkRoabs{} cm, 캐나다에서 \ndrefCablkRoabs{} cm 줄였고, 학습에서 제외한 지역에서도 Stefan 모델 수준을 유지하였으며, 레나 델타와 외부 CALM 지역에서 지역 라벨 5--25개로 오차를 줄였다. 직접 ML은 라벨을 모두 쓴 지역 안 검증에서는 점추정 오차가 더 낮았으나, 레나 델타 라벨 없이 학습하면 그 지역에서 오차가 Stefan 모델보다 \ndrefLeregDzabs{} cm 컸다. 같은 셀 · 같은 해에서 이 모델은 제품이 학습하지 않은 지역에서 공개 제품 5종보다 정확하였고, 제품이 인근 관측으로 학습한 지역에서는 Wei 등 · Liu 등과 같은 수준이었다. 따라서 Stefan 모델을 기준으로 두고 잔차만 학습하는 구조는 대상 지역 라벨이 적거나 없을 때 가장 유용하며, 이는 영구동토 지역의 일반적인 관측 조건이다. 이 워크플로는 적용 가능 영역과 예측 구간을 갖춘 연도별 1 km ALT 지도와, 다른 영구동토 지역에서 ALT 모델과 제품을 평가하고 보정하는 절차를 제공한다.
